\documentclass[11pt,a4paper]{article}

% ── Encoding & Language ──
\usepackage[T1]{fontenc}
\usepackage[utf8]{inputenc}
\usepackage[nil]{babel}

% ── Fonts (matching original: Palatino body, Helvetica sans, Courier mono) ──
\usepackage{mathpazo}        % Palatino (URWPalladioL equivalent)
\usepackage[scaled=0.90]{helvet}   % Helvetica (BeraSans equivalent)
\usepackage{courier}  % Courier (BeraSansMono equivalent)

% ── Layout ──
\usepackage[a4paper, top=2.5cm, bottom=2.5cm, left=2.8cm, right=2.8cm]{geometry}
\usepackage{microtype}
\usepackage{parskip}
\setlength{\parindent}{0pt}
\setlength{\parskip}{4pt}

% ── Colors ──
\usepackage[dvipsnames,table]{xcolor}
\definecolor{darkred}{HTML}{8C1414}
\definecolor{brownday}{HTML}{681900}
\definecolor{headergray}{HTML}{7F807F}
\definecolor{tablegray}{HTML}{F2F2F2}
\definecolor{titlegray}{HTML}{595959}

% ── Tables ──
\usepackage{booktabs}
\usepackage{array}
\usepackage{longtable}
\usepackage{tabularx}
\usepackage{colortbl}
\newcolumntype{L}[1]{>{\raggedright\arraybackslash}p{#1}}
\newcolumntype{C}[1]{>{\centering\arraybackslash}p{#1}}

% ── Headers & Footers ──
\usepackage{fancyhdr}
\pagestyle{fancy}
\fancyhf{}
\renewcommand{\headrulewidth}{0pt}
\fancyhead[L]{\footnotesize\sffamily\color{headergray}\leftmark}
\fancyfoot[C]{\small\thepage\ | \pageref{lastpage}}

% ── Section formatting (matching original: BeraSans-Bold, darkred, 17.8pt) ──
\usepackage{titlesec}
\titleformat{\section}
  {\sffamily\bfseries\Large\color{darkred}}
  {}{0em}{}
\titlespacing{\section}{0pt}{18pt}{8pt}

\titleformat{\subsection}
  {\sffamily\bfseries\large\color{brownday}}
  {}{0em}{}
\titlespacing{\subsection}{0pt}{12pt}{4pt}

\titleformat{\subsubsection}
  {\rmfamily\bfseries\normalsize}
  {}{0em}{}
\titlespacing{\subsubsection}{0pt}{8pt}{3pt}

% ── Mark for header ──
\usepackage{nameref}
\renewcommand{\sectionmark}[1]{\markboth{#1}{}}

% ── TOC ──
\usepackage{tocloft}
\renewcommand{\cftsecfont}{\sffamily\bfseries}
\renewcommand{\cftsecpagefont}{\sffamily\bfseries\color{darkred}}
\renewcommand{\cftsubsecfont}{\rmfamily}
\renewcommand{\cftsubsecpagefont}{\rmfamily\color{darkred}}
\renewcommand{\cftsecleader}{\cftdotfill{\cftdotsep}}
\setlength{\cftbeforesecskip}{4pt}

% ── Hyperref ──
\usepackage[hidelinks,colorlinks=false]{hyperref}

% ── Misc ──
\usepackage{enumitem}
\setlist[itemize]{leftmargin=1.5em, itemsep=1pt, parsep=0pt, topsep=4pt}
\renewcommand{\labelitemi}{$\bullet$}

\newcommand{\notetitle}[1]{\subsubsection*{#1}}

% ── Cover page ──
\newcommand{\makecover}{%
  \thispagestyle{empty}
  \vspace*{3cm}
  {\fontsize{42}{46}\selectfont\rmfamily\bfseries\color{darkred}LabBook\par}
  \vspace{0.3cm}
  {\large\sffamily\bfseries Diario del Proyecto\par}
  \vspace{0.2cm}
  {\large\rmfamily\color{titlegray}Análisis de Yahoo! Answers con XML y eXist-db\par}
  \vspace{0.1cm}
  {\large\rmfamily\color{titlegray}YahooAnswers-XML\par}
  \vspace{1.5cm}
  {\normalsize\rmfamily\bfseries Matteo Amagliani $\cdot$ Manuel Avilés Rodríguez\par}
  \vspace{0.8cm}
  {\normalsize\rmfamily Proyecto: YahooAnswers-XML --- Análisis con Bases de Datos Documentales\par}
  {\normalsize\rmfamily Asignatura: Bases de Datos Documentales --- Módulo XML\par}
  \vspace{0.3cm}
  {\large\rmfamily\bfseries\color{darkred}Marzo 2026\par}
  \newpage
}

% ── Activity table ──
\newcommand{\activitytable}[2]{%
  \begin{center}
  \begin{tabular}{C{2.5cm} L{7.8cm} C{2.8cm} C{1.8cm}}
  \toprule
  \rowcolor{white}
  \textbf{Hora} & \textbf{Actividad} & \textbf{Responsable} & \textbf{Duración} \\
  \midrule
  #1
  \bottomrule
  \end{tabular}
  \end{center}
  \vspace{-4pt}
  {\small\textbf{Cuadro:} \textit{#2}}
  \vspace{6pt}
}

\begin{document}

\makecover

% ════════ TOC ════════
\tableofcontents
\newpage

% ════════════════════════════════════════════════════════
\section{Semana 1: Preparación (16--17 Marzo 2026)}
\markboth{Semana 1: Preparación (16--17 Marzo 2026)}{}

\subsection*{1.\quad Planificación}

\subsection{Domingo 16 de marzo}

\activitytable{%
  10:00--11:00 & Configuración repositorio: crear repo GitHub con estructura de carpetas, README, .gitignore & Ambos & 1 h \\
  \rowcolor{tablegray}
  11:00--12:00 & Exploración XML: analizar estructura de los 4 ficheros XML extraídos del corpus & Ambos & 1 h \\
}{Registro de actividades --- Domingo 16 Mar.}

\notetitle{Notas:}
\begin{itemize}
  \item Se creó el repositorio \textit{YahooAnswers-XML} en GitHub con la estructura de carpetas definida en la planificación: \texttt{src/}, \texttt{existdb/}, \texttt{xquery/}, \texttt{xsd/}, \texttt{xslt/}, \texttt{data/}.
  \item Se añadió el fichero \texttt{.gitignore} para excluir los XML de datos (demasiado grandes para el repositorio) y los entornos virtuales de Python.
  \item Se exploraron los 4 ficheros XML ya extraídos del corpus original (12 GB): \textit{preguntas\_politica\_y\_gobierno.xml} (12.976 docs), \textit{preguntas\_noticias\_y\_actualidad.xml} (674 docs), \textit{preguntas\_politica\_e\_governo.xml} (2.118 docs) y \textit{preguntas\_notizie\_ed\_eventi.xml} (1.184 docs).
  \item Se verificó que la estructura XML de todos los ficheros es consistente: elemento raíz \texttt{<preguntas>} con atributo \texttt{categoria}, y elementos \texttt{<document>} con los campos del corpus original.
\end{itemize}

\subsection{Lunes 17 de marzo}

\activitytable{%
  10:00--11:30 & Instalación eXist-db: instalar y configurar eXist-db localmente. Probar panel de administración y REST API. & Ambos & 1,5 h \\
}{Registro de actividades --- Lunes 17 Mar.}

\notetitle{Notas:}
\begin{itemize}
  \item Se descargó el instalador de eXist-db 6.x desde GitHub Releases y se ejecutó mediante \texttt{java -jar eXist-db-setup-6.x.x.jar}.
  \item Se verificó el acceso al panel de administración en \texttt{http://localhost:8080/exist/} y al editor de consultas eXide en \texttt{http://localhost:8080/exist/apps/eXide/}.
  \item Se probó la REST API con una petición básica y se confirmó que responde correctamente con HTTP 200.
\end{itemize}

\notetitle{Dificultades encontradas:}
\begin{itemize}
  \item La instalación de eXist-db requirió Java 8+ preinstalado. En el sistema había Java 17, que resultó compatible sin problemas adicionales.
  \item Al primer arranque, el proceso tardó aproximadamente 45 segundos en estar completamente operativo, lo que inicialmente generó confusión al intentar acceder al panel antes de tiempo.
\end{itemize}

\newpage

% ════════════════════════════════════════════════════════
\section{Semana 1 (cont.): Esquema XSD (17--19 Marzo 2026)}
\markboth{Semana 1 (cont.): Esquema XSD (17--19 Marzo 2026)}{}

\subsection{Martes 17 de marzo (tarde)}

\activitytable{%
  14:00--16:00 & XSD original: construir XSD que describe el formato actual de los XML del corpus & Manuel Avilés & 2 h \\
}{Registro de actividades --- Martes 17 Mar.}

\notetitle{Notas:}
\begin{itemize}
  \item Se construyó \textit{yahoo\_answers\_original.xsd} con la estructura real del corpus: elemento raíz \texttt{<preguntas>} con atributo obligatorio \texttt{categoria}, y tipo complejo \texttt{DocumentType} con todos los campos del corpus.
  \item Campos obligatorios: \texttt{uri}, \texttt{subject}, \texttt{bestanswer}, \texttt{nbestanswers}, \texttt{cat}, \texttt{maincat}, \texttt{subcat}, \texttt{date}, \texttt{res\_date}, \texttt{lastanswerts}, \texttt{qlang}, \texttt{qintl}, \texttt{language}, \texttt{id}, \texttt{best\_id}.
  \item Campos opcionales (\texttt{minOccurs="0"}): \texttt{content} (ausente en el 31,3\% de los documentos) y \texttt{vot\_date} (ausente en el 53,2\%).
  \item Se definieron dos tipos simples restringidos: \texttt{CodigoIdioma} (patrón \texttt{[a-z]\{2\}}) y \texttt{UserId} (patrón \texttt{u[0-9]+}).
\end{itemize}

\subsection{Miércoles 18 de marzo}

\activitytable{%
  10:00--12:30 & XSD ampliado: diseñar extensión del esquema con métricas de calidad y nodos de red social & Matteo Amagliani & 2,5 h \\
}{Registro de actividades --- Miércoles 18 Mar.}

\notetitle{Notas:}
\begin{itemize}
  \item Se diseñó \textit{yahoo\_answers\_extended.xsd} como extensión del esquema original, añadiendo dos nuevos elementos opcionales al tipo \texttt{DocumentType}:
  \begin{itemize}
    \item \texttt{<metricas\_calidad>}: campos derivados como \texttt{tiene\_content}, \texttt{longitud\_subject}, \texttt{longitud\_bestanswer}, \texttt{num\_respuestas}, \texttt{fue\_votado} y \texttt{tiempo\_resolucion\_seg}.
    \item \texttt{<red\_social>}: anotaciones para el análisis de redes, con campos \texttt{autor\_es\_respondedor}, \texttt{subcats\_usuario} y \texttt{peso\_enlace}.
  \end{itemize}
  \item Todos los campos nuevos son opcionales para mantener compatibilidad con los XML existentes del corpus.
\end{itemize}

\notetitle{Dificultades encontradas:}
\begin{itemize}
  \item Decidir el nivel de extensión del esquema: se optó por añadir solo campos que aportan valor analítico real y que pueden calcularse a partir de los datos existentes, en lugar de crear campos especulativos.
\end{itemize}

\subsection{Jueves 19 de marzo (mañana)}

\activitytable{%
  10:00--12:00 & Validación XML: probar los XSD contra los 4 ficheros XML con \texttt{validate\_xml.py} & Manuel Avilés & 2 h \\
}{Registro de actividades --- Jueves 19 Mar. (mañana)}

\notetitle{Notas:}
\begin{itemize}
  \item Se ejecutó \texttt{validate\_xml.py} sobre los 4 ficheros XML usando \texttt{lxml}. Los 4 pasaron la validación contra el XSD original sin errores.
  \item Se confirmó que el campo \texttt{content} aparece en algunos documentos pero no en otros, validando el \texttt{minOccurs="0"} del esquema.
\end{itemize}

\newpage

% ════════════════════════════════════════════════════════
\section{Semana 2: XSLT y Carga en eXist-db (19--20 Marzo 2026)}
\markboth{Semana 2: XSLT y Carga en eXist-db (19--20 Marzo 2026)}{}

\subsection{Jueves 19 de marzo (tarde)}

\activitytable{%
  14:00--16:00 & XSLT Inmigración/Actualidad: transformación HTML navegable para las categorías en español & Manuel Avilés & 2 h \\
  \rowcolor{tablegray}
  14:00--16:00 & XSLT Immigrazione/Attualità: transformación HTML navegable para las categorías en italiano & Matteo Amagliani & 2 h \\
}{Registro de actividades --- Jueves 19 Mar. (tarde)}

\notetitle{Notas:}
\begin{itemize}
  \item Se crearon dos ficheros XSLT: \textit{es\_html.xsl} (cabecera azul, textos en español) e \textit{it\_html.xsl} (cabecera roja, textos en italiano).
  \item Ambas transformaciones generan una página HTML con: cabecera con nombre de categoría, total de preguntas, y una tarjeta por documento con el \texttt{subject}, el \texttt{bestanswer} resaltado en verde y un desplegable para las demás respuestas.
  \item Se aplicaron localmente con \texttt{lxml}: la transformación del fichero más grande (\textit{politica\_y\_gobierno.xml}, 12.976 documentos) generó un HTML de aproximadamente 45 MB en 8 segundos.
\end{itemize}

\notetitle{Decisiones tomadas:}
\begin{itemize}
  \item Se usó XSLT 1.0 para garantizar compatibilidad con \texttt{lxml} sin necesidad de Saxon. Las funciones avanzadas de XSLT 2.0 no eran necesarias para esta transformación.
\end{itemize}

\subsection{Viernes 20 de marzo (mañana)}

\activitytable{%
  10:00--11:30 & Carga XML: cargar los 4 XML en eXist-db usando REST API (script \texttt{cargar\_existdb.py}) & Ambos & 1,5 h \\
}{Registro de actividades --- Viernes 20 Mar. (mañana)}

\notetitle{Notas:}
\begin{itemize}
  \item Se ejecutó \texttt{cargar\_existdb.py}, que crea las colecciones \texttt{/db/yahooanswers/es/} e \texttt{/db/yahooanswers/it/} y sube cada fichero XML mediante HTTP PUT.
  \item Tiempos de carga: \textit{noticias\_y\_actualidad.xml} (1 s), \textit{notizie\_ed\_eventi.xml} (3 s), \textit{politica\_e\_governo.xml} (5 s), \textit{politica\_y\_gobierno.xml} (38 s).
  \item Se verificó la carga ejecutando \texttt{count(collection("/db/yahooanswers")//document)} en eXide: resultado 16.952, coincide con el total esperado.
\end{itemize}

\notetitle{Dificultades encontradas:}
\begin{itemize}
  \item El fichero \textit{politica\_y\_gobierno.xml} ($\approx$95 MB) fue rechazado con HTTP 413 en el primer intento. Se solucionó editando \texttt{\$EXIST\_HOME/etc/jetty/webdefault.xml} aumentando el límite de tamaño de petición a 200 MB.
  \item Se aumentó el timeout del script de carga a 300 segundos para el fichero más grande.
\end{itemize}

\newpage

% ════════════════════════════════════════════════════════
\section{Semana 3: Análisis con XQuery (20--22 Marzo 2026)}
\markboth{Semana 3: Análisis con XQuery (20--22 Marzo 2026)}{}

\subsection{Viernes 20 de marzo (tarde)}

\activitytable{%
  14:00--16:00 & Consultas básicas XQuery: estadísticas generales, distribución por subcategorías, campos ausentes & Matteo Amagliani & 2 h \\
}{Registro de actividades --- Viernes 20 Mar. (tarde)}

\notetitle{Notas:}
\begin{itemize}
  \item Se desarrolló \textit{01\_estadisticas.xq} con conteos por categoría, media de respuestas, distribución de subcategorías y porcentaje de documentos sin campo \texttt{content}.
  \item Resultados: total 16.952 documentos, 161.365 respuestas. \textit{Política y Gobierno} (ES) concentra el 76,5\% de las preguntas. La media de respuestas por pregunta varía de 7,03 (Noticias y Actualidad) a 12,02 (Politica e governo IT).
  \item La consulta se ejecutó en 0 ms según el tiempo devuelto por eXist-db, gracias a sus índices internos de colección.
\end{itemize}

\subsection{Sábado 21 de marzo}

\activitytable{%
  10:00--12:30 & Análisis Redes Sociales: XQuery sobre usuarios top, simetría pregunta/respuesta, clustering por categoría & Manuel Avilés & 2,5 h \\
  \rowcolor{tablegray}
  14:00--16:00 & Calidad de Datos: XQuery sobre nulos, longitudes, fechas inconsistentes y anomalías & Matteo Amagliani & 2 h \\
}{Registro de actividades --- Sábado 21 Mar.}

\notetitle{Notas:}
\begin{itemize}
  \item Se desarrolló \textit{02\_redes\_sociales.xq}. Resultados más relevantes:
  \begin{itemize}
    \item De los 12.881 usuarios distintos, el 47,9\% solo pregunta, el 36,7\% solo responde y únicamente el 15,4\% hace ambas cosas. Baja simetría consistente con la literatura.
    \item El usuario más activo como respondedor (\texttt{u1012406}) acumula 398 mejores respuestas, concentradas en 2 categorías únicas (clustering elevado).
    \item Solo 2 usuarios participan activamente tanto en la comunidad hispanohablante como en la italohablante, lo que confirma que ambas comunidades son prácticamente independientes.
  \end{itemize}
  \item Se desarrolló \textit{03\_calidad\_datos.xq}. Resultados más relevantes:
  \begin{itemize}
    \item El campo \texttt{content} está ausente en el 31,3\% de los documentos y \texttt{vot\_date} en el 53,2\%.
    \item Longitud media del \texttt{bestanswer}: 438,5 caracteres. El máximo es 3.996 caracteres y el mínimo 1 carácter.
    \item 677 documentos tienen una mejor respuesta de menos de 20 caracteres (anomalía de calidad).
    \item No se detectaron inconsistencias de fechas (ningún documento con \texttt{date} $>$ \texttt{res\_date}).
    \item El 16,8\% de las preguntas se resolvieron el mismo día de publicación.
  \end{itemize}
\end{itemize}

\notetitle{Dificultades encontradas:}
\begin{itemize}
  \item La consulta de redes sociales tardó 14 segundos en ejecutarse sobre los 16.952 documentos. Aunque aceptable, se identificó como un área de mejora con índices adicionales en eXist-db.
\end{itemize}

\subsection{Domingo 22 de marzo}

\activitytable{%
  10:00--12:00 & Consultas avanzadas: XQuery comparativa ES vs IT, palabras clave en política, preguntas sin respuesta & Ambos & 2 h \\
}{Registro de actividades --- Domingo 22 Mar.}

\notetitle{Notas:}
\begin{itemize}
  \item Se desarrolló \textit{04\_consultas\_avanzadas.xq} con tres secciones: comparativa ES vs IT, palabras clave políticas y preguntas sin respuesta.
  \item Comparativa ES vs IT: los usuarios italianos responden de media 11,32 veces por pregunta frente a 9,08 de los hispanohablantes, y sus respuestas son ligeramente más largas (443,8 vs 437,3 caracteres de media).
  \item Palabras clave: \textit{presidente} aparece en 732 preguntas de política, seguido de \textit{gobierno} (429) y \textit{ley} (273). El término italiano \textit{governo} aparece en 112 preguntas.
  \item Preguntas sin respuesta: 0. Todas las preguntas del corpus tienen al menos una respuesta registrada.
\end{itemize}

\notetitle{Dificultades encontradas:}
\begin{itemize}
  \item La primera versión de \textit{04\_consultas\_avanzadas.xq} incluía ordenación de los 17.000 documentos y agrupación con \texttt{group by}. Ambas operaciones causaron timeout de 300 segundos en eXist-db al no disponer de índices específicos. Se simplificó la consulta eliminando esas secciones.
\end{itemize}

\newpage

% ════════════════════════════════════════════════════════
\section{Semana 4: Documentación y Entrega (22--25 Marzo 2026)}
\markboth{Semana 4: Documentación y Entrega (22--25 Marzo 2026)}{}

\subsection{Lunes 23 de marzo}

\activitytable{%
  10:00--14:00 & Redacción memoria T1: arquitectura, esquemas, XSLT, consultas, resultados y capturas & Ambos & 4 h \\
}{Registro de actividades --- Lunes 23 Mar.}

\notetitle{Notas:}
\begin{itemize}
  \item Se redactaron las secciones principales de la memoria: introducción al corpus, descripción de la estructura XML, diseño del XSD original y extendido, explicación de las transformaciones XSLT, descripción de la carga en eXist-db y análisis de resultados de las 4 consultas XQuery.
  \item Se incluyeron capturas de pantalla del editor eXide con los resultados de cada consulta, y ejemplos de las páginas HTML generadas por los XSLT.
\end{itemize}

\subsection{Martes 24 de marzo}

\activitytable{%
  10:00--12:00 & Dashboard de resultados: desarrollo del visor HTML interactivo con gráficos Chart.js & Matteo Amagliani & 2 h \\
  \rowcolor{tablegray}
  12:00--13:00 & Colchón de desviaciones: corrección de errores menores en XSD y ajustes de consultas & Ambos & 1 h \\
  14:00--15:30 & Revisión cruzada de la memoria y correcciones finales & Ambos & 1,5 h \\
}{Registro de actividades --- Martes 24 Mar.}

\notetitle{Notas:}
\begin{itemize}
  \item Se desarrolló \texttt{ver\_resultados.py}, un script que genera un dashboard HTML interactivo con 4 pestañas (Estadísticas, Redes Sociales, Calidad de Datos, Consultas Avanzadas) y gráficos Chart.js a partir de los XML de resultados generados por eXist-db.
  \item El dashboard usa un diseño oscuro con tipografías Inter y JetBrains Mono, y una paleta de colores personalizada (cian para ES, naranja-salmón para IT).
  \item Se realizó revisión cruzada de la memoria.
\end{itemize}

\subsection{Miércoles 25 de marzo}

\activitytable{%
  10:00--11:00 & Revisión y empaquetado T1: verificar repo, nombrado final de ficheros, subida a GitHub & Ambos & 1 h \\
  \rowcolor{tablegray}
  11:00--12:00 & Preparación de la presentación oral para la defensa del 26 de marzo & Ambos & 1 h \\
}{Registro de actividades --- Miércoles 25 Mar.}

\notetitle{Notas:}
\begin{itemize}
  \item Se verificó la estructura completa del repositorio: todos los ficheros en su carpeta correcta, README actualizado con el flujo completo paso a paso.
  \item Se exportaron los documentos entregables: memoria en PDF, LabBook, documento de contribuciones.
  \item Entrega realizada antes de las 20:00 del miércoles 25 de marzo, dentro del plazo establecido.
\end{itemize}

\notetitle{Documentos generados:}
\begin{itemize}
  \item \textit{2\_Proyecto\_AmaglianiAviles.pdf}: Memoria completa del proyecto.
  \item \textit{1\_LabBook\_AmaglianiAviles.pdf}: LabBook actualizado.
  \item \textit{0\_Contribuciones\_AmaglianiAviles.pdf}: Documento de contribuciones con firmas.
\end{itemize}

\newpage

% ════════════════════════════════════════════════════════
\section*{Resumen de Actividades (al 25 de marzo)}
\addcontentsline{toc}{section}{Resumen de Actividades (al 25 de marzo)}
\markboth{Resumen de Actividades (al 25 de marzo)}{}

\begin{center}
\begin{tabular}{L{11cm} C{3.5cm}}
\toprule
\textbf{Actividad} & \textbf{Estado} \\
\midrule
Configuración del repositorio & Completado \\
\rowcolor{tablegray}
Instalación y configuración de eXist-db & Completado \\
Exploración de los 4 ficheros XML del corpus & Completado \\
\rowcolor{tablegray}
XSD original (formato actual del corpus) & Completado \\
XSD ampliado (métricas de calidad y red social) & Completado \\
\rowcolor{tablegray}
Validación XML contra XSD con \texttt{lxml} & Completado \\
XSLT Inmigración/Actualidad (ES $\to$ HTML) & Completado \\
\rowcolor{tablegray}
XSLT Immigrazione/Attualità (IT $\to$ HTML) & Completado \\
Carga de los 4 XML en eXist-db (REST API) & Completado \\
\rowcolor{tablegray}
\textit{01\_estadisticas.xq}: estadísticas generales & Completado \\
\textit{02\_redes\_sociales.xq}: simetría y usuarios top & Completado \\
\rowcolor{tablegray}
\textit{03\_calidad\_datos.xq}: completitud, anomalías, fechas & Completado \\
\textit{04\_consultas\_avanzadas.xq}: comparativa ES/IT, palabras clave & Completado \\
\rowcolor{tablegray}
Dashboard HTML interactivo (\texttt{ver\_resultados.py}) & Completado \\
Memoria del proyecto & Completado \\
\rowcolor{tablegray}
LabBook & Completado \\
Documento de contribuciones & Completado \\
\rowcolor{tablegray}
Empaquetado y entrega final & Completado \\
\bottomrule
\end{tabular}
\end{center}

\newpage

% ════════════════════════════════════════════════════════
\section*{Anexo: Resumen de Tiempo Dedicado y Dificultades}
\addcontentsline{toc}{section}{Anexo: Resumen de Tiempo Dedicado y Dificultades}
\markboth{Anexo: Resumen de Tiempo Dedicado y Dificultades}{}

\subsection*{A.1\quad Tiempo Dedicado}

\notetitle{Por día}

\begin{center}
\begin{tabular}{L{3cm} C{1.8cm} L{9cm}}
\toprule
\textbf{Fecha} & \textbf{Horas} & \textbf{Actividades principales} \\
\midrule
Dom 16 mar & 2,0 h & Configuración repositorio, exploración XML \\
\rowcolor{tablegray}
Lun 17 mar & 2,5 h & Instalación eXist-db, XSD original \\
Mar 17 mar (tarde) & 2,5 h & XSD original (cont.), inicio XSD ampliado \\
\rowcolor{tablegray}
Mié 18 mar & 2,5 h & XSD ampliado, inicio XSLT \\
Jue 19 mar & 5,5 h & XSLT ES e IT, validación XML, carga eXist-db \\
\rowcolor{tablegray}
Vie 20 mar & 4,5 h & Carga eXist-db, consultas básicas XQuery \\
Sáb 21 mar & 4,5 h & Análisis Redes Sociales, Calidad de Datos \\
\rowcolor{tablegray}
Dom 22 mar & 4,0 h & Consultas avanzadas, corrección timeout \\
Lun 23 mar & 4,0 h & Redacción memoria T1 \\
\rowcolor{tablegray}
Mar 24 mar & 4,5 h & Dashboard, colchón, revisión cruzada \\
Mié 25 mar & 2,0 h & Empaquetado y entrega final \\
\midrule
\rowcolor{tablegray}
\textbf{Total} & \textbf{38,5 h} & \\
\bottomrule
\end{tabular}
\end{center}

\vspace{8pt}
\notetitle{Por componente del equipo}

\begin{center}
\begin{tabular}{L{8cm} C{2.5cm} C{2.5cm}}
\toprule
\textbf{Componente} & \textbf{Horas} & \textbf{Porcentaje} \\
\midrule
Matteo Amagliani & 19,0 h & 49\% \\
\rowcolor{tablegray}
Manuel Avilés Rodríguez & 19,5 h & 51\% \\
\midrule
\rowcolor{tablegray}
\textbf{Total} & \textbf{38,5 h} & \textbf{100\%} \\
\bottomrule
\end{tabular}
\end{center}

\vspace{8pt}
\notetitle{Por tipo de actividad}

\begin{center}
\begin{tabular}{L{9cm} C{2cm} C{2.5cm}}
\toprule
\textbf{Tipo de actividad} & \textbf{Horas} & \textbf{Porcentaje} \\
\midrule
Preparación (repositorio, exploración XML) & 3,0 h & 8\% \\
\rowcolor{tablegray}
Instalación y configuración eXist-db & 2,5 h & 6\% \\
Diseño XSD original y ampliado & 6,5 h & 17\% \\
\rowcolor{tablegray}
Transformaciones XSLT & 4,0 h & 10\% \\
Carga en eXist-db & 1,5 h & 4\% \\
\rowcolor{tablegray}
Consultas XQuery (4 consultas) & 8,5 h & 22\% \\
Dashboard de resultados & 2,0 h & 5\% \\
\rowcolor{tablegray}
Documentación y memoria & 8,0 h & 21\% \\
Empaquetado y entrega & 2,5 h & 7\% \\
\midrule
\rowcolor{tablegray}
\textbf{Total} & \textbf{38,5 h} & \textbf{100\%} \\
\bottomrule
\end{tabular}
\end{center}

\vspace{12pt}
\subsection*{A.2\quad Dificultades Encontradas}

\notetitle{Dificultad 1: Tiempo de arranque de eXist-db}
\textbf{Tipo:} Técnica\\
\textbf{Descripción:} Al primer arranque de eXist-db, el proceso tardó aproximadamente 45 segundos en estar completamente operativo. Intentar acceder al panel antes de que terminase de arrancar devolvía errores de conexión.\\
\textbf{Solución:} Esperar a que el log de arranque indicase \texttt{Server has started, listening on: http://localhost:8080/} antes de hacer ninguna petición.\\
\textbf{Impacto:} Bajo --- solo afectó al primer uso.

\vspace{6pt}
\notetitle{Dificultad 2: Límite de tamaño de petición en eXist-db}
\textbf{Tipo:} Técnica\\
\textbf{Descripción:} El fichero \textit{politica\_y\_gobierno.xml} ($\approx$95 MB) fue rechazado con HTTP 413 en el primer intento de carga, ya que el servidor Jetty embebido en eXist-db tiene un límite de tamaño de petición por defecto.\\
\textbf{Solución:} Se editó el fichero \texttt{\$EXIST\_HOME/etc/jetty/webdefault.xml} aumentando el parámetro \texttt{maxFormContentSize} a 200 MB. También se aumentó el timeout del script de carga a 300 segundos.\\
\textbf{Tiempo perdido:} $\sim$30 minutos.

\vspace{6pt}
\notetitle{Dificultad 3: Timeout en consulta XQuery con ordenación global}
\textbf{Tipo:} Técnica\\
\textbf{Descripción:} La primera versión de \textit{04\_consultas\_avanzadas.xq} incluía la ordenación de los 17.000 documentos por número de respuestas para obtener el top 3 por categoría. Esta operación causaba timeout a los 300 segundos sin devolver ningún resultado.\\
\textbf{Solución:} Se eliminó la sección de top preguntas más respondidas de la consulta 4. Se mantuvo únicamente la comparativa ES/IT, las palabras clave y las preguntas sin respuesta, que son operaciones de agregación simple sin ordenación global.\\
\textbf{Tiempo perdido:} $\sim$1 hora (dos intentos fallidos y rediseño).

\vspace{6pt}
\notetitle{Dificultad 4: Compatibilidad XSLT 1.0 vs 2.0}
\textbf{Tipo:} Técnica\\
\textbf{Descripción:} La plantilla XSLT inicialmente usaba funciones de XPath 2.0. Aunque eXist-db soporta XSLT 2.0, la librería \texttt{lxml} de Python solo soporta XSLT 1.0, por lo que la transformación local fallaba.\\
\textbf{Solución:} Se reescribieron los XSLT usando únicamente funciones de XSLT 1.0 compatibles con \texttt{lxml}.\\
\textbf{Tiempo perdido:} $\sim$20 minutos.

\vspace{6pt}
\notetitle{Dificultad 5: Campos opcionales en el XSD}
\textbf{Tipo:} Diseño\\
\textbf{Descripción:} El campo \texttt{content} está ausente en el 31,3\% de los documentos y \texttt{vot\_date} en el 53,2\%. Definir estos campos como obligatorios hacía fallar la validación de casi un tercio del corpus.\\
\textbf{Solución:} Se marcaron \texttt{content} y \texttt{vot\_date} con \texttt{minOccurs="0"} en el XSD, reflejando correctamente la naturaleza opcional de estos campos en el corpus real.\\
\textbf{Impacto:} Bajo --- detectada rápidamente en la primera ejecución de \texttt{validate\_xml.py}.

\vspace{20pt}
\hrule
\vspace{6pt}
{\small\textit{Última actualización: 25 de marzo de 2026}}

\label{lastpage}
\end{document}
